\documentclass[11pt,a4paper]{ctexart}
\usepackage[margin=2.4cm]{geometry}
\usepackage{amsmath,amssymb}
\usepackage{fancyhdr}
\usepackage{enumitem}
\linespread{1.18}
\setlength{\headheight}{15pt}
\setlength{\parskip}{0pt}

\pagestyle{fancy}
\fancyhf{}
\fancyhead[L]{\small Team 2405424}
\fancyhead[R]{\small \thepage}
\renewcommand{\headrulewidth}{0.4pt}
\fancypagestyle{plain}{\fancyhf{}\renewcommand{\headrulewidth}{0pt}}

\ctexset{
  section/format = \Large\bfseries\raggedright,
  subsection/format = \large\bfseries\raggedright,
  subsubsection/format = \normalsize\bfseries\raggedright
}
\setcounter{secnumdepth}{3}

\begin{document}

%====================== 摘要页 ======================
\thispagestyle{empty}
\begin{center}
{\small 2024 MCM/ICM\quad 选题 A\quad 队伍控制号 2405424\quad 摘要页（Summary Sheet）}\\[18pt]
{\LARGE\bfseries 七鳃鳗：中庸之道的生动例证}\\[8pt]
{\large （Lamprey: a Vivid Example of the Doctrine of the Mean）}\\[10pt]
\end{center}

\begin{center}
{\large\bfseries 摘\quad 要}
\end{center}

\vspace{6pt}
作为地球上最古老的生物之一，七鳃鳗的生存机制始终是研究者关注的话题。当讨论七鳃鳗与资源可获得性时，其种群能够改变性别比这一现象便成为无法回避的议题。研究已经发现，资源的可获得性会影响七鳃鳗种群的性别比。

为了研究物种依据资源可获得性改变性别比的能力，我们以 Lotka-Volterra 捕食者—猎物模型与 Nicholson-Bailey 寄主—寄生物模型作为建模基础，进而构建了理想参考模型、恒定性别比模型与性别比变化模型，并成功求解了该问题。

针对任务 1，我们首先发现：随着七鳃鳗或其他寄生物的引入，生态系统中的资源总量出现下降，而存活物种数在含七鳃鳗的模型中最高。由此我们得出结论：与普通寄生物相比，能够改变性别比的七鳃鳗对其他生物造成的生存压力更小，因此有利于提高生态系统的多样性。

针对任务 2，我们人为地使环境资源恶化，并分别绘制了恒定性别比模型与性别比变化模型中七鳃鳗数量的变化曲线。显而易见，常规寄生类物种的种群规模受到环境的严重影响，而能够改变性别比的七鳃鳗则非常稳定，这是七鳃鳗最大的种群优势。然而有得必有失：对环境风险具有更强韧性的代价是更低的种群规模，这正是七鳃鳗种群的劣势所在。

针对任务 3，我们收集了大量数据，通过熵权法分析了三种模型最终的生态系统稳定性，并得出结论：寄生物的引入导致生态系统稳定性下降，但七鳃鳗能够改变性别比这一特性显著缓解了这一现象。

针对任务 4，我们随后整合了这三个模型，以探究在生态系统中引入七鳃鳗对其他常规寄生物的影响。在恶化条件下，我们发现：在缺乏七鳃鳗的生态系统中，寄生物抵御外部环境的能力较弱。然而在引入七鳃鳗之后，其种群规模呈现出稳定的周期性波动。这放大了生态系统的负反馈调节作用，也使原有寄生物对不利环境的抵抗力显著增强。

此后，为检验模型，我们首先从多个方面讨论了实验结果的合理性。接着，我们从人类活动的扰动与种群初值的变化两个角度对模型进行了灵敏度分析，发现模型对不确定性因素具有更强的抗干扰能力，并且适用于不同的种群初值，这恰好体现了模型的可靠性与普适性。最后，我们分析了模型的优缺点，并提出了相应的优化方案。

\vspace{6pt}
\noindent\textbf{关键词：}性别比变化，理想参考模型，恒定性别比模型，性别比变化模型，熵权法

\newpage
\pagestyle{fancy}

%%====================== 正文 ======================
\section{引言}

\subsection{研究背景}

\subsubsection{七鳃鳗的生活习性}

七鳃鳗从受精卵到成体的整个生命周期通常需要 5—8 年。受精卵孵化成蠕虫状幼体，栖息于沙质溪流底部的小洞穴中。幼体以碎屑和小型植物藻类为食，经过三至六年后变为寄生性成体。成体迁移至湖泊和海洋之中，以鱼类为食。作为典型的洄游生物，七鳃鳗有部分时间生活在海中。它们在秋季进入大河，在河流下游越冬，并于次年 5 月至 6 月水温达到约 15 摄氏度时返回上游产卵。雌雄亲体在产卵后都会死亡。

寄生是七鳃鳗主要的生存方式。成体七鳃鳗通过带齿的吸盘状口吸附在寄主鱼类身上，吸食其血液。这种寄生行为不仅损害寄主鱼类的健康、降低其生存与繁殖能力，还会影响鱼类种群的数量与结构，从而影响水生生态系统的平衡。

\subsubsection{七鳃鳗的种群动态}

由于七鳃鳗在某些文化中被视为美味，近年来人类的大量捕捞已导致其种群数量下降。在考虑七鳃鳗的捕食、被捕食与寄生关系时，性别比的变化也可能对其生态角色产生显著影响。例如，若雄性比例上升导致繁殖率下降，则七鳃鳗对寄主鱼类的影响可能减弱，从而影响相关的食物链。此外，性别比的变化还可能影响七鳃鳗的遗传多样性与适应能力，进而影响其长期存续能力。

\subsection{问题重述}

\begin{itemize}
\item 针对任务 1，我们需要基于七鳃鳗能够改变性别比这一基本特性建立数学模型，该模型应能反映七鳃鳗改变性别比的特性对生态系统的影响，并依据所建模型得出相应结论。
\item 针对任务 2，我们需要运用并改进上述模型，从模型结果中导出七鳃鳗种群的优点与缺点。
\item 针对任务 3，我们需要建立生态系统稳定性的评价标准，并通过比较引入七鳃鳗前后该指标的变化，确定并验证七鳃鳗能够改变性别比这一特性对生态系统稳定性的作用。
\item 针对任务 4，我们需要进一步细化模型，找出是否存在能够从含有七鳃鳗种群的生态系统中获利的物种（例如寄生物），并阐述其获利的机制。
\end{itemize}

\subsection{文献综述}

由于七鳃鳗与生态系统中的其他物种之间存在种间竞争、寄生等关系，为更好地描述七鳃鳗在生态系统中的地位与作用，我们需要了解关于种间关系的已有数学模型。据文献 [3]、[4]，现有能够描述种间关系的模型主要是描述捕食关系的 Lotka-Volterra 捕食者—猎物模型，以及描述寄生关系的 Nicholson-Bailey 寄主—寄生物模型。在此基础上，我们建立模型的思路如图 2 所示。

如图 2 可见，由于所有物种都从环境资源中获取能量，因此可以把环境资源视为猎物、把每个物种视为捕食者，并可用 Lotka-Volterra 捕食者—猎物模型来描述物种与环境资源之间的关系。又由于多个物种捕食同一环境资源，当用数学模型描述物种与环境资源之间的关系时，各物种之间的种间竞争关系也被间接地描述了。

\subsection{我们的工作}

\begin{enumerate}
\item 在 Nicholson-Bailey 寄主—寄生物模型与 Lotka-Volterra 捕食者—猎物模型的基础上，我们将建立理想参考模型、恒定性别比模型与性别比变化模型。
\item 通过调整各模型中的生态系统参数，我们将模拟不同气候条件下不同物种的生存状况以及资源总量的变化。
\item 我们将分析所得数据并比较上述三个模型，讨论七鳃鳗改变性别比的能力对七鳃鳗自身种群的优缺点，以及对生态系统存活物种数、资源总量和生态系统稳定性等因素的影响。
\item 结合这三个模型，我们将通过对恒定性别比模型与性别比变化模型作合理改进，讨论七鳃鳗对性别比为 $1:1$ 的常规寄生物的影响。
\item 最后，我们将同时通过改变各物种种群规模的初值以及加入人类捕捞因素的扰动这两种方式，对模型进行灵敏度分析。
\end{enumerate}

\section{模型准备}

\subsection{假设}

\begin{itemize}
\item \textbf{假设 1：}不考虑除人类以外的寄生类物种的天敌。即群落中物种与寄生生物之间的关系以寄生与种间竞争为主。

理由：据文献 [5]，七鳃鳗等寄生类物种在北美五大湖属于入侵物种，说明其天敌极少，可以忽略。

\item \textbf{假设 2：}种群中雌雄七鳃鳗的生存能力大致相同。

理由：由文献 [6] 可知，雌雄七鳃鳗的生存能力差异很小；该假设可避免因雌性或雄性生存能力不足而导致某一性别大量死亡的极端情形，从而保证模型预测结果的可靠性。

\item \textbf{假设 3：}生态系统中的所有物种都能直接从环境资源中获取能量。

理由：简单分析可知，当该假设成立时，用模型 1 反映不同物种之间的竞争关系更为有效，同时也在一定程度上简化了模型。
\end{itemize}

以上是主要假设，其他假设将在文中适当位置给出。

\subsection{符号说明}

为更好地描述模型，我们赋予下列符号实际含义：

\begin{itemize}
\item $t$：七鳃鳗的繁殖轮数
\item $N(t)$：$t$ 时刻七鳃鳗的种群规模
\item $s(t)$：七鳃鳗平均每个个体所占的资源份额
\item $m(s)$：平均资源份额为 $s$ 时七鳃鳗中雄性个体的比例
\item $\lambda(t)$：$t$ 时刻的资源总量
\item $s_0$：能够保证七鳃鳗种群正常生存的平均资源最小值
\item $N_i(t)$：$t$ 时刻群落中第 $i$ 个物种的种群规模
\item $n$：生态系统物种数（不含七鳃鳗）的初值
\item $n_1$：生态系统物种数（不含七鳃鳗）的终值
\item $ESI$：生态系统稳定性指数
\end{itemize}

其余符号将在文中适当位置列出。
\section{模型的建立}

为了更有说服力地反映七鳃鳗改变性别比对生态系统的影响，我们基于三种情景建立了三类模型，分别对应三类生态系统：不含七鳃鳗的理想参考生态系统、含性别比恒为 $2:1$ 的常规寄生生物的生态系统，以及含正常七鳃鳗的生态系统。第一个生态系统作为后续评价生态系统相关指标的参照；第二与第三个生态系统可以相互比较，以反映七鳃鳗改变性别比的功能对生态系统的影响。

\subsection{理想参考模型}

在该模型所代表的生态系统中不存在七鳃鳗。为简化模型，可以假设这一理想生态系统中的所有主要物种彼此之间均为种间竞争关系。由于这些物种共享同一环境资源，因此可以假设环境资源为猎物、这些物种为捕食者。参照 Lotka-Volterra 捕食者—猎物模型 [3]，我们可以提出一个更完整的、用于描述该生态系统的模型。理想参考生态系统的基本结构如图 3 所示。

首先，以环境资源为研究对象，可列出其差分方程如下：
\begin{equation}
\lambda(t+1)-\lambda(t)=\sigma\lambda(t)-q\lambda^{2}(t)-\sum_{i=1}^{n}\mu_{i}\lambda(t)N_{i}(t)
\end{equation}
其中 $\lambda(t)$ 表示 $t$ 时刻的环境资源总量，$\sigma$ 表示种群的自然增长率；$\lambda^{2}(t)$ 项表示构成环境资源的基础生物（如植物微生物）的内部竞争，$\lambda(t)N_{i}(t)$ 项表示群落中第 $i$ 个物种对该资源的摄取速率。$q$ 为底层生物内部竞争的强度，$\mu_{i}$ 为第 $i$ 个物种对应的环境需求强度系数，二者均与物种自身的性质有关。

然后以第 $i$ 个物种为研究对象，同样容易列出其差分方程：
\begin{equation}
N_{i}(t+1)-N_{i}(t)=-a_{i}N_{i}(t)+b_{i}N_{i}(t)\lambda(t)-p_{i}N_{i}^{2}(t)\quad (i=1,2,\cdots\cdots,n)
\end{equation}
其中 $N_{i}(t)$ 表示 $t$ 时刻第 $i$ 个物种的总数量，$a_{i}$ 表示第 $i$ 个物种在缺乏食物时的自然衰减率。与前一方程相同，$N_{i}(t)\lambda(t)$ 项表示群落中第 $i$ 个物种对资源的摄取速率，$N_{i}^{2}(t)$ 项表示第 $i$ 个物种的种内竞争。$b_{i}$ 与 $p_{i}$ 分别为相应种间或种内关系的竞争强度系数，与不同物种自身的习性有关。

将两类差分方程联立，得到方程组：
\begin{equation}
\begin{cases}
N_{i}(t+1)-N_{i}(t)=-a_{i}N_{i}(t)+b_{i}N_{i}(t)\lambda(t)-p_{i}N_{i}^{2}(t)\quad (i=1,2,\cdots\cdots,n)\\[6pt]
\lambda(t+1)-\lambda(t)=\sigma(t)\lambda(t)-q\lambda^{2}(t)-\sum\limits_{i=1}^{n}\mu_{i}\lambda(t)N_{i}(t)
\end{cases}
\end{equation}

由于 $q$ 是种内竞争强度，受外部不确定性影响很小，故可认为 $q$ 的取值是常数。而对于 $q$ 以外的系数，由于物种的变异性与外部条件的不确定性，我们将其视为随机变量。其中，$\sigma(t)$ 是关于繁殖轮数 $t$ 的随机变量，它通过表示资源总量受外部环境变化影响的不确定性来反映生态系统内部的不确定性；$a_{i}$、$b_{i}$、$p_{i}$ 与 $\mu_{i}$ 都是关于 $i$ 的随机变量，表示不同物种之间竞争能力的差异。经过计算机软件试验，可以确定相关系数的适当取值。

我们假设初始物种类型数为 10，每个物种的初始总数量为 $240{,}000$，初始环境资源为 $4.5\times10^{9}$ kJ，由此可得到各物种种群规模随繁殖轮数的变化，具体情形如图 4 所示。

在图 4 中，红线表示环境资源总量，蓝线表示不同物种。由图 4 以及多次实验计算出的平均值，我们可以得出结论：理想生态系统的环境资源约为 $4.89\times10^{9}$ kJ，其所能容纳的主要物种数大致为 8.88。在本次模拟中有两个物种处于优势地位，一个物种处于劣势地位，另有两个物种灭绝。这可以作为衡量七鳃鳗对生态系统影响的参照。

\subsection{恒定性别比模型（基于含性别比恒为 $2:1$ 的常规寄生类物种的生态系统）}

据文献 [7]，七鳃鳗的性别比一般最高为 $5:1$、最低为 $1:1$，即雄性比例落在区间 $[0.50,0.83]$ 之内。由于本模型的存在是为了反映七鳃鳗改变性别比这一能力的作用，因此可以假设本模型中常规寄生类物种的性别比恒为 $2:1$，即雄性个体比例为 0.67。因此，该模型主要基于性别比恒为 $2:1$ 的常规寄生类物种、多个寄主以及多个不会被寄生的常规物种。

就常规寄生类物种种群规模的变化而言，由于雌雄比例保持不变，我们考虑以下因素：个体繁殖的自然增长率、个体的平均资源分配，以及与其他物种的种间关系。这些因素都与该物种的总量相联系，它们作用于该物种的总量，同时也受其反作用。它们的相互关系如图 5 所示。

如图 5 所示，种群总量主要受三个因素影响：自然增长（出生率与死亡率）、环境约束以及与其他物种的关系，下文将分别讨论并建模。

\subsubsection{自然增长}

该寄生类物种的种群必然因繁殖而出现自然增长。由于繁殖产生的后代数量高度依赖于雌性数量，并考虑到存在部分个体因非繁殖因素在繁殖前死亡，可以假设后代数量与成功繁殖的雌性数量成正比，用方程表示如下：
\begin{equation}
N(t+1)=k(1-m)^{\alpha}N(t)
\end{equation}
其中 $N(t+1)$ 表示繁殖产生的后代数量，$(1-m)$ 表示雌性比例，$\alpha$ 表示比例修正系数，本文经计算机模拟后取 $\alpha=1.25$ 时模型更为合理；$k$ 为自然增长系数，与雌性的平均产仔数以及因非繁殖因素导致的种群死亡率有关。

\subsubsection{环境约束}

环境对物种的约束一般体现在平均资源占有量的大小上。我们用 $s$ 表示 $t$ 时刻的平均资源占有量。由平均资源占有量的定义，$s$ 的公式如下：
\begin{equation}
s(t)=\frac{\text{资源总量}}{\text{种群总量}}=\frac{\lambda(t)}{N(t)+\sum\limits_{i=1}^{n}N_{i}(t)}
\end{equation}
其中 $\lambda(t)$ 为 $t$ 时刻的资源总量，$N(t)+\sum_{i=1}^{n}N_{i}(t)$ 为 $t$ 时刻的种群总量。然而经计算机软件模拟，当性别比恒定时，平均资源占有量 $s$ 一般变化不大，对种群的影响较小，在此类模型中可暂时忽略；但该因素将在下一个模型中发挥重要作用。

\subsubsection{与其他物种的关系}

根据假设 1，寄生类物种与其他物种的关系主要分为种间竞争与寄生两类，其中寄生是该物种与其他物种之间的主要种间关系。下面分别讨论。

首先，对于种间竞争，我们可以参照理想模型中对种间竞争的数学描述。竞争由多个物种捕食同一物种而产生。假设群落中考虑在内的物种总数（不含该寄生类物种）为 $n$，其中 $z$ 个物种可被该寄生类物种寄生，则仅考虑种间竞争的相关方程组如下：
\begin{equation}
\begin{cases}
N_{i}(t+1)-N_{i}(t)=-a_{i}N_{i}(t)+b_{i}N_{i}(t)\lambda(t)-p_{i}N_{i}^{2}(t)\quad (i=1,2,\cdots\cdots,n-z)\\[6pt]
\lambda(t+1)-\lambda(t)=\sigma(t)\lambda(t)-q\lambda^{2}(t)-\beta\lambda(t)N(t)-\sum\limits_{i=1}^{n}\mu_{i}\lambda(t)N_{i}(t)
\end{cases}
\end{equation}
相对于理想模型，第二个方程中多出的 $\beta\lambda(t)N(t)$ 项表示寄生类物种对环境资源的消耗。对于可被寄生类物种寄生的那 $z$ 个物种，它们之间的关系更为复杂，将在后文单独讨论。

其次，对于群落各物种之间的寄生关系，主要的寄生结构是寄生类物种作为寄生物、另外 $z$ 个物种作为寄主。为研究它们之间的关系，需要先求解单对寄生物与其寄主之间的关系。

1935 年，爱尔兰学者 A.J. Nicholson 与英国学者 V.A. Bailey 提出了寄主—寄生物模型，形式如下 [4]：
\begin{equation}
\begin{cases}
Q(t+1)=rQ(t)\mathrm{e}^{-aP(t)}\\
P(t+1)=Q(t)\left[1-\mathrm{e}^{-aP(t)}\right]
\end{cases}
\end{equation}
其中 $Q$ 表示寄主数量，$P$ 表示寄生物数量，$r$ 表示寄主未被寄生时的自然增长率，$a$ 表示寄生物的搜索概率。然而该模型在自然环境中波动很大，与实际情形并不十分吻合，因此后来的学者对其进行了改进，得到如下方程组 [4]：
\begin{equation}
\begin{cases}
Q(t+1)=Q(t)\mathrm{e}^{\rho\left(1-Q(t)\right)-\frac{P(t)}{1+aQ(t)}}\\[6pt]
P(t+1)=\phi Q(t)\left(1-\mathrm{e}^{-\frac{P(t)}{1+aQ(t)}}\right)
\end{cases}
\end{equation}
其中 $Q$ 与 $P$ 分别表示寄主的种群密度与寄生物的种群密度，$\rho$、$a$ 与 $\phi$ 为相关系数。由于上述 $z$ 个物种不仅与寄生类物种存在寄生关系，它们彼此之间还存在这种竞争关系，因此需要将上述种间竞争模型与寄生模型整合，从而导出这 $z$ 个物种的种群数量变化方程，如下：
\begin{equation}
\begin{cases}
N_{i1}(t+1)=(1-a_{i})N_{i}(t)+b_{i}N_{i}(t)\lambda(t)-p_{i}N_{i}^{2}(t)\\[6pt]
N_{i2}(t+1)=N_{i}(t)\mathrm{e}^{\rho\left[1-\frac{N_{i}(t)}{10^{6}}\right]-\frac{N(t)}{10^{6}+a_{i}N_{i}(t)}}\\[6pt]
N_{i}(t+1)=\gamma N_{i1}(t+1)+(1-\gamma)N_{i2}(t+1)\\[6pt]
\qquad\qquad i=n-z+1,\ n-z+2,\ \cdots\cdots,\ n
\end{cases}
\end{equation}
其中 $N_{i1}(t+1)$ 表示寄主的种间竞争部分，$N_{i2}(t+1)$ 表示寄主的被寄生部分。方程组中的 $10^{6}$ 是为把物种数量换算为种群密度而选取的合适面积大小；$\gamma$ 是寄主与寄生物之间的种间竞争关系相对于寄生关系的强度系数。

最后，我们综合上述三个因素，并将涉及寄生类物种种群数量的相关方程整合，得到如下关于寄生类物种种群变化的方程：
\begin{equation}
N(t+1)=k(1-m)^{1.25}\left(1-\mathrm{e}^{-\left(\frac{N(t)}{10^{6}}+\sum\limits_{i=n-z+1}^{n}a_{i}N_{i}(t)\right)}\right)\sum_{i=n-z+1}^{n}\phi_{i}N_{i}(t)
\end{equation}
对于该方程，由于种群的天然增长与寄生均被认为会直接影响寄生类物种的种群规模，我们将式（4）与改进的 Nicholson-Bailey 寄主—寄生物模型适当整合得到该方程，用于刻画该寄生类物种种群数量的变化。

综上分析，我们可以导出性别比不变时描述整个生态系统的总方程组，如下：
\begin{equation}
\begin{cases}
N(t+1)=k(1-m)^{1.25}\left(1-\mathrm{e}^{-\left(\frac{N(t)}{10^{6}}+\sum\limits_{i=n-z+1}^{n}a_{i}N_{i}(t)\right)}\right)\sum\limits_{i=n-z+1}^{n}\phi_{i}N_{i}(t)\\[6pt]
\lambda(t+1)-\lambda(t)=\sigma(t)\lambda(t)-q\lambda^{2}(t)-\beta\lambda(t)N(t)-\sum\limits_{i=1}^{n}\mu_{i}\lambda(t)N_{i}(t)\\[6pt]
N_{i}(t+1)-N_{i}(t)=-a_{i}N_{i}(t)+b_{i}N_{i}(t)\lambda(t)-p_{i}N_{i}^{2}(t)\\[4pt]
\qquad\qquad (i=1,2,\cdots\cdots,n-z)\\[6pt]
N_{i}(t+1)=\gamma\left[(1-a_{i})N_{i}(t)+b_{i}N_{i}(t)\lambda(t)-p_{i}N_{i}^{2}(t)\right]\\[4pt]
\qquad\qquad +(1-\gamma)\left[N_{i}(t)\mathrm{e}^{\rho\left(1-\frac{N_{i}(t)}{10^{6}}\right)-\frac{N(t)}{10^{6}+a_{i}N_{i}(t)}}\right]\\[4pt]
\qquad\qquad (i=n-z+1,\ n-z+2,\ \cdots\cdots,\ n)
\end{cases}
\end{equation}
其中我们默认 $m=0.67$，即雌雄比为 $2:1$；该方程组中涉及环境因素与物种竞争差异的相关系数，我们都将其设为适当的随机变量。在后文中，通过将该模型与性别比变化模型进行比较，我们可以得出改变性别比对环境的影响，以及这一特性给寄生类物种种群带来的优势与劣势。

\subsection{性别比变化模型（基于含正常七鳃鳗的生态系统）}

该模型关注七鳃鳗对生态系统的影响，因此对应的是主要物种组成为七鳃鳗、七鳃鳗的寄主以及不被七鳃鳗寄生的常规物种的生态系统。根据题意，七鳃鳗个体的性别受外部环境条件影响。当资源稀缺、环境恶劣时，七鳃鳗个体的性别更可能是雄性；当资源丰富、条件良好时，个体的性别更可能是雌性。在本模型中，我们用平均资源量来反映环境条件。

基于性别比随平均资源占有量变化的总趋势，可以假设当平均资源占有量接近 0 时，雄性比例趋于接近 1；而当资源十分充足时，雌雄比趋于接近 $1:1$ [7]。不过，由于资源过少时种群无法生存，因此当平均资源占有量为 $s_{0}$ 时雄性比例最大；由文献 [7] 可知此时雌雄比约为 $1:5$，即雄性比例为 0.83。因此可以认为，该关系在一定程度上类似于 Logistic 模型，于是参照 Logistic 模型，可以给出 $m$ 与 $s$ 之间的近似关系方程如下：
\begin{equation}
m(s)=\frac{1}{2}+\frac{1}{2+\mathrm{e}^{b(s-s_{0})}}\quad (s\ge s_{0})
\end{equation}
其中 $b$ 为该曲线的伸缩因子，用于调整出合适的曲线形状。该曲线的图形表示如图 6 所示。

由图 6 可见，该曲线大致呈现类似字母 S 的形状，与实际情形相符；图中实线所指示的部分是模型主要涉及的部分（即雄性份额 $m$ 落在区间 $[0.50,0.83]$ 内的情形）。

经过上述讨论可以得出，雄性比例 $m$ 与平均资源占有量 $s$ 之间存在显著关系。由于种群的天然增长率与性别比密切相关，进而影响七鳃鳗的种群总量 $N$，而种群总量 $N$ 又再次作用于平均资源占有量 $s$。因此可以引入生态系统内部的相互关系，如图 7 所示。

其中 $+$ 号表示箭头首尾两端因素之间为正相关，$-$ 号表示箭头首尾两端因素之间为负相关。图 7 中所表示的多个因素之间的关系构成多个闭合回路，通过负反馈调节体现生态系统的稳定性。

我们将性别比的变化作为一个重要影响因素纳入恒定性别比模型的方程组，得到如下综合模型：
\begin{equation}
\begin{cases}
N(t+1)=k(1-m)^{1.25}\left(1-\mathrm{e}^{-\left(\frac{N(t)}{10^{6}}+\sum\limits_{i=n-z+1}^{n}a_{i}N_{i}(t)\right)}\right)\sum\limits_{i=n-z+1}^{n}\phi_{i}N_{i}(t)\\[6pt]
s(t)=\dfrac{\lambda(t)}{N(t)+\sum\limits_{i=1}^{n}N_{i}(t)}\\[6pt]
m(s)=\dfrac{1}{2}+\dfrac{1}{2+\mathrm{e}^{b(s-s_{0})}}\\[6pt]
\lambda(t+1)-\lambda(t)=\sigma(t)\lambda(t)-q\lambda^{2}(t)-\beta\lambda(t)N(t)-\sum\limits_{i=1}^{n}\mu_{i}\lambda(t)N_{i}(t)\\[6pt]
N_{i}(t+1)-N_{i}(t)=-a_{i}N_{i}(t)+b_{i}N_{i}(t)\lambda(t)-p_{i}N_{i}^{2}(t)\\
\qquad\qquad (i=1,2,\cdots\cdots,n-z)\\[6pt]
N_{i}(t+1)=\gamma\left[(1-a_{i})N_{i}(t)+b_{i}N_{i}(t)\lambda(t)-p_{i}N_{i}^{2}(t)\right]\\
\qquad\qquad +(1-\gamma)\left[N_{i}(t)\mathrm{e}^{\rho\left(1-\frac{N_{i}(t)}{10^{6}}\right)-\frac{N(t)}{10^{6}+a_{i}N_{i}(t)}}\right]\\
\qquad\qquad (i=n-z+1,\ n-z+2,\ \cdots\cdots,\ n)
\end{cases}
\end{equation}
该模型由恒定性别比模型改进而来，因此与原模型具有很强的相似性。主要改进在于性别比 $m$：为符合实际，性别比 $m$ 不再是一个数值，而是依据式（12）与其他因素相关联的变量。该模型成功刻画了性别比的变化，为后文讨论改变性别比的作用奠定了基础。

\subsection{生态系统的评价指标与稳定性分析}

据文献 [8]，为评估模型对生态系统稳定性的影响，需要考虑两个主要评估指标：生态资源总量与生物多样性。我们选择资源总量作为评估指标，主要因为生态资源总量直接影响各群落的生理活力，并决定生态系统的结构、功能与生产力。此外，大量研究发现生物多样性水平与生态系统稳定性密切相关，生物多样性越高，相应生态系统的稳定性越好。因此，生物多样性也是评估生态系统稳定性的关键指标之一，本文用生态系统最终存活的物种数来表示。

对于存在七鳃鳗的模型，为综合上述两个指标来评价生态系统的稳定性，我们采用熵权法求出不同指标的权重。具体步骤在后文说明。
\section{问题求解}

为直观展示模型，我们假设生态系统中的初始物种数 $n$ 为 10，其中有三种生物是七鳃鳗的寄主。

\subsection{七鳃鳗性别比变化对生态系统的影响}

此处我们关注七鳃鳗改变性别比对生态系统资源总量、其他物种种群规模以及存活物种数的影响。为研究生态系统中的资源总量与存活物种数，我们进行了多次模拟实验，将这些实验随机且均匀地分为若干组，取每一部分实验的平均值，得到若干组数据，具体数值与实验组序号之间的关系如图 8 与图 9 所示。

由图 8 与图 9 可见，经多次试验，不同生态系统中生物多样性与生态资源总量之间的大小关系相对稳定。容易得出结论：与理想生态系统相比，引入主要依赖寄生生活的相应常规物种会提高生态系统的物种多样性，并降低环境资源总量。然而，当引入七鳃鳗时，与其它常规寄生类物种相比，其改变性别比的能力提高了外部生态系统的物种多样性，而环境资源总量并未发生明显变化。为解释这一现象并考察对其他物种种群规模的影响，我们对引入常规寄生类物种与引入七鳃鳗两种情形进行了模拟，并以繁殖轮数 $t$ 为横轴、以不同物种的数量与环境资源总量为纵轴绘图表示生态系统状况，如图 10 与图 11 所示。

根据图 10 与图 11 的信息，引入常规寄生类物种给生态系统中其他物种带来的压力明显大于引入七鳃鳗所带来的压力，这体现为图 11 中其他物种的物种数总体低于图 10 中的情形。正因如此，相对于常规物种，七鳃鳗促进了生态系统中的物种多样性，因为七鳃鳗对其他物种的压力小于常规物种的压力。从引入常规物种到引入七鳃鳗，寄生生物数量减少而其他生物数量增加，因此对环境资源总量的影响并不显著。

因此，与理想生态系统相比，引入能够改变性别比的七鳃鳗可以提高生态系统的物种多样性，并降低生态系统的资源总量。此外，与引入常规寄生类物种相比，引入七鳃鳗能够进一步提高生态系统的多样性，并缓解其他物种之间的竞争压力。

\subsection{改变性别比对七鳃鳗的优势与劣势}

经过多次计算机模拟，我们发现，为使改变性别比对种群的作用更为显著，需要不断恶化环境条件。通过放大环境资源条件中的不确定性，即增大方程组（11）中第二个方程的随机变量 $\sigma(t)$ 的方差，我们使环境条件变化更为剧烈，从而体现环境的恶化。

为直观呈现结果，我们用计算机绘制了不同模型对应的图形，以繁殖轮数 $t$ 为横轴、以不同物种的数量与环境资源总量为纵轴。于是，七鳃鳗性别比恒定的模型的实验结果如图 12 所示。

由图 12 可见，随着繁殖轮数增加，环境资源的波动显著增大，代表环境恶化。除此之外，图 12 还显示七鳃鳗的种群规模也随环境资源的剧烈振荡而显著振荡，说明若性别比恒定，七鳃鳗种群受环境变化的影响较大，难以保持种群规模稳定，不利于种群的长期存续。若考虑七鳃鳗改变性别比的能力，通过计算机模拟性别比变化模型可得结果，如图 13 所示。

根据该模型，尽管环境资源的波动仍然很大，七鳃鳗的种群规模却基本保持稳定不变，说明其对环境具有更强的抵抗力，更适合种群的长期存续，这或许正是七鳃鳗能够在地球上存活 3.6 亿年的原因！

再次回顾上述两图，经比较可见，虽然改变性别比的能力使七鳃鳗种群对外部扰动更具抵抗力，但正因这一能力，七鳃鳗的种群规模仅维持在 $1\times10^{6}$ 左右。另一方面，当性别比恒定时，七鳃鳗的最小种群规模约为 $2\times10^{6}$，这表明改变性别比的能力以提高种群抵御外部扰动能力的代价是降低了种群规模。

因此我们可以得出结论：改变性别比的能力给七鳃鳗种群带来的优势是种群稳定性提高，给该物种带来的劣势是种群规模受到限制。

\subsection{改变性别比对七鳃鳗所在生态系统稳定性的影响}

由上可知，生态系统的稳定性主要受存活物种数与环境资源总量控制。因此，可以认为生态系统稳定性的计算公式如下：
\begin{equation}
ESI=\omega_{1}\lambda+\omega_{2}n_{1}
\end{equation}
其中 $ESI$ 表示生态系统稳定性，$\lambda$ 与 $n_{1}$ 分别表示经多次实验得到的生态系统资源总量与存活物种数的平均值，$\omega_{1}$ 与 $\omega_{2}$ 为相应指标的权重。

为求得 $ESI$ 的值，应分别求出各指标的权重以及生态资源总量与存活物种数这两个指标的平均值。对于权重，我们采用熵权法根据已知数据求得。由熵权法的定义可知：首先，应通过计算机模拟收集理想参考模型、恒定性别比模型与性别比变化模型若干组数据的生态资源总量与存活物种数，并分别进行标准化；然后，利用标准化后的数据分别计算生态资源总量的信息熵 $E_{1}$ 与存活物种数的信息熵 $E_{2}$，具体公式如下：
\begin{equation}
E_{j}=-\frac{\sum\limits_{i=1}^{n_{0}}p_{ij}\ln p_{ij}}{\ln n_{0}}\quad (j=1,2)
\end{equation}
其中 $p_{ij}$ 表示该类中各数据所占的比例，$n_{0}$ 表示样本总量。此后，可求得两类中生态资源总量与存活物种数各自的信息冗余度，进而计算出两类中生态资源总量与生物多样性的最终权重。最后，计算相应模型的生态资源总量与存活物种数的均值，并将其代入方程，即可得到三者的 $ESI$ 值；通过比较其大小便可得出七鳃鳗对生态系统稳定性的影响。

具体评估结果为：理想参考模型的存活物种数为 0.75、资源量为 0.55、稳定性为 0.63；恒定性别比模型的存活物种数为 0.77、资源量为 0.38、稳定性为 0.54；性别比变化模型的存活物种数为 0.84、资源量为 0.45、稳定性为 0.61。

根据上述数据可以得出结论：引入常规寄生类物种会使生态系统稳定性下降，而七鳃鳗改变性别比的能力能够显著缓解这一现象。

\subsection{改变性别比对七鳃鳗相对于其他物种的优势}

根据题目提示，为弄清七鳃鳗能否为其他物种（例如寄生物）带来优势，我们考察了七鳃鳗种群对另一种寄生物的影响。为突出这一作用，我们同样通过增大随机变量 $\sigma(t)$ 的方差来适当调整环境资源变化的幅度，以体现更为严酷的环境资源状况。作为参照，首先应导出种群中不存在七鳃鳗时寄生物种群规模随繁殖轮数的变化。由于大多数寄生类物种的性别比为 $1:1$，我们在恒定性别比模型的基础上将雄性比例由 $2/3$ 改为 $1/2$，模拟结果如图 14 所示。

由图可见，当环境资源变化较为剧烈时，寄生物种群数量的变化与环境资源的变化高度重合，寄生物的曲线也呈现出不规则的大幅振荡。这一现象表明，在缺乏七鳃鳗的情况下，寄生物种群受环境影响较大，其抵御不利环境的能力较弱。

在引入七鳃鳗后，生态系统中存在两类被寄生的物种。为便于计算机模拟，我们假设七鳃鳗寄生三个物种，而寄生物寄生另外两个物种，此外还有四个不会被寄生的物种。由于存在两类寄生类物种，我们需要将恒定性别比模型与性别比变化模型整合，得到的具体方程组为：
\begin{equation}
\begin{cases}
N(t+1)=k(1-m)^{1.25}\left(1-\mathrm{e}^{-\left(\frac{N(t)}{10^{6}}+\sum\limits_{i=5}^{7}a_{i}N_{i}(t)\right)}\right)\sum\limits_{i=5}^{7}\phi_{i}N_{i}(t)\\[6pt]
s(t)=\dfrac{\lambda(t)}{N(t)+\sum\limits_{i=1}^{n}N_{i}(t)}\\[6pt]
m(s)=\dfrac{1}{2}+\dfrac{1}{2+\mathrm{e}^{b(s-s_{0})}}\\[6pt]
N_{10}(t+1)=k(1-m_{10})^{1.25}\left(1-\mathrm{e}^{-\left(\frac{N_{10}(t)}{10^{6}}+\sum\limits_{i=8}^{9}a_{i}N_{i}(t)\right)}\right)\sum\limits_{i=8}^{9}\phi_{i}N_{i}(t)\quad \left(m_{10}=\dfrac{1}{2}\right)\\[6pt]
\lambda(t+1)-\lambda(t)=\sigma(t)\lambda(t)-q\lambda^{2}(t)-\beta\lambda(t)N(t)-\sum\limits_{i=1}^{10}\mu_{i}\lambda(t)N_{i}(t)\\[6pt]
N_{i}(t+1)-N_{i}(t)=-a_{i}N_{i}(t)+b_{i}N_{i}(t)\lambda(t)-p_{i}N_{i}^{2}(t)\quad (i=1,2,3,4)
\end{cases}
\end{equation}
\begin{equation}
\begin{cases}
N_{i}(t+1)=\gamma\left[(1-a_{i})N_{i}(t)+b_{i}N_{i}(t)\lambda(t)-p_{i}N_{i}^{2}(t)\right]\\
\qquad\qquad +(1-\gamma)\left[N_{i}(t)\mathrm{e}^{\rho\left(1-\frac{N_{i}(t)}{10^{6}}\right)-\frac{N(t)}{10^{6}+a_{i}N_{i}(t)}}\right]\quad (i=5,6,7)\\[6pt]
N_{i}(t+1)=\gamma\left[(1-a_{i})N_{i}(t)+b_{i}N_{i}(t)\lambda(t)-p_{i}N_{i}^{2}(t)\right]\\
\qquad\qquad +(1-\gamma)\left[N_{i}(t)\mathrm{e}^{\rho\left(1-\frac{N_{i}(t)}{10^{6}}\right)-\frac{N_{10}(t)}{10^{6}+a_{i}N_{i}(t)}}\right]\quad (i=8,9)
\end{cases}
\end{equation}

其中 $N_{i}(t)\,(i=1,2,3,4)$ 表示四个不会被寄生的物种各自的种群数量，$N_{i}(t)\,(i=5,6,7)$ 表示将被七鳃鳗寄生的三个寄主的种群数量，$N_{i}(t)\,(i=8,9)$ 表示将被寄生物寄生的两个寄主的种群数量，$N_{10}(t)$ 表示寄生物的种群数量，$N(t)$ 表示七鳃鳗的种群数量，$m_{10}$ 表示寄生物的性别比，此处设为 $1/2$。通过计算机模拟可得结果如图 15 所示。

由图 15 可见，当七鳃鳗数量达到某一数值后，寄生物数量与七鳃鳗数量呈现周期性波动，且两者种群数量曲线的峰值位置与对方的谷值位置相对应，表明二者之间存在协作关系、相互控制，使种群数量呈现稳定的周期波动。因此，通过比较引入七鳃鳗前后寄生物种群规模的变化可以发现：当环境条件剧烈波动时，引入七鳃鳗后寄生物种群数量的变化曲线不再是大尺度不规则的振荡曲线，而是具有固定周期与振幅的振荡曲线，说明其对外部环境的抵抗力得到了增强。

综上讨论可以得出结论：引入能够改变性别比的七鳃鳗，可通过与寄生物形成协作关系来增强其抵御外部环境的能力，使寄生物在严酷的环境变化面前表现出更好的生存优势。

\section{模型的合理性分析与灵敏度检验}

\subsection{合理性分析}

经我们判断，该模型有三个与实际情形高度吻合的特点，如下：

\begin{itemize}
\item 在该模型中，各物种的数量在经过一定时间后都会趋于平稳，这与生物学中生态系统具有一定稳定性的结论相一致。
\item 模型中七鳃鳗种群的波动与环境资源总量的波动几乎实现了同频同向的振动，体现了物种对环境的依赖性。
\item 该模型成功预测并清晰展示了七鳃鳗与寄生物及其他物种之间可能存在的相互作用。它们之间的相互制约导致动态平衡，使其种群数量平稳增减，体现了生态系统具有一定自我调节能力的生态学原理。
\end{itemize}

这三点与实际情形非常吻合，从多个方面说明了模型的合理性。

\subsection{灵敏度检验}

\subsubsection{人类活动}

在某些地区，七鳃鳗是重要的食物来源，可被人类捕捞，从而导致七鳃鳗的种群规模骤降。为检验模型的合理性，我们考虑人类在 $T_{0}$ 时刻对七鳃鳗进行捕捞，则捕捞前后该种群的关系如图 16 所示。

由图 16 可见，人类捕捞之后生态系统并未崩溃，而是迅速恢复到人类捕捞之前的状态，说明模型对人类等外部因素的扰动不敏感，体现了模型的可靠性与稳定性。

\subsubsection{种群初值}

由于不同生态系统中各物种种群规模的初值不同，为检验模型的普适性，我们不断改变模型中主要物种种群规模的初值，研究生态系统中最终存活物种数的均值、环境资源总量与初始种群规模之间的关系，结果如图 17 所示。

由图 17 可见，随着种群规模初值的增大，该生态系统中最终存活的主要物种总数稳定在 9.2 左右；但由于种群数量增加，对环境的消耗也随之增大，可见环境资源总量持续下降。这些现象非常合理，因此模型能够针对不同的初始种群规模给出非常恰当的结果，体现了模型的普适性。

经过上述两个方面，可以得出结论：我们的模型具有良好的抵御外部不确定性的能力与优良的普适性。对于不同尺度、包含不确定性的生态系统，模型的分析与预测是可靠的。

\section{模型的优缺点}

\subsection{优点}

\begin{itemize}
\item 模型具有灵活性：通过调整模型中的参数，可以模拟不同的生态系统，从而得到贫瘠或丰富的环境资源，或处于优势或劣势的种群。
\item 模型具有真实性：模型中的参数可以调整以模拟自然灾害，从而判断外部环境动荡对生态系统的影响。
\item 七鳃鳗寄生物种的数量可以自由调整，以模拟七鳃鳗在不同群落中的生存状况，有利于对七鳃鳗种群进行深入研究。
\end{itemize}

\subsection{缺点}

\begin{itemize}
\item 计算机模拟模型中物种总数（不含七鳃鳗）仅为 10，数量过少，偏离实际生态系统。

可能的解决方案：增加模型中的物种数与方程数量。

\item 模型中的食物网不够复杂，生态系统结构较为简单。

可能的解决方案：引入高级消费者，同时依据不同营养级对模型进行分层。
\end{itemize}

\section{结论}

本文基于 Lotka-Volterra 捕食者—猎物模型与 Nicholson-Bailey 寄主—寄生物模型建立了三个模型，并通过合理比较这三个模型的运行结果，针对不同问题得出了不同结论。结合我们的结论可以得出：寄生类物种的引入会给其他物种带来压力并降低生态系统的稳定性，但七鳃鳗改变性别比的独特功能不仅缓解了这些负面影响，还增强了种群中其他寄生生物的稳定性，这足以体现七鳃鳗对生态系统的中庸之道。

\section*{参考文献}

\begin{enumerate}
\item[{[1]}] W. Defen, ``Fishery conservation in the great lakes of north america: Invasion and control of seven gilled eels in the sea of the great lakes of north america,'' \emph{Chinese Fisheries}, vol. 000, 2012.
\item[{[2]}] https://cn.bing.com/images/search.
\item[{[3]}] Z. Zhenliang, W. Runxin, L. Liyun, and Y. Xiangqin, ``The influence of fear effect to the lotka-volterra predator-prey system with predator has other food resource,'' \emph{ADVANCES IN DIFFERENCE EQUATIONS}, vol. 2020, 2020.
\item[{[4]}] W. ruiwu, ``Nicholson-bailey host-parasite model,'' \emph{Encyclopedia of China}, 2023.
\item[{[5]}] M. J. Siefkes, ``Use of physiological knowledge to control the invasive sea lamprey (petromyzon marinus) in the laurentian great lakes,'' \emph{CONSERVATION PHYSIOLOGY}, vol. 5, 2017.
\item[{[6]}] L. Peng, C. Hui, and Z. Wenge, ``Studies on the winter hermaphroditic morphology and individual fecundity of the japanese lamprey,'' \emph{Freshwater Fisheries}, 2008.
\item[{[7]}] P. S. Maitland, ``Ecology of the river, brook and sea lamprey,'' \emph{Conserving Natura 2000 Rivers Ecology series No.5}, 2003.
\item[{[8]}] H. Nianpeng, X. Li, and H. Honglin, ``The methods of evaluation ecosystem quality: Ideal reference and key parameters,'' \emph{Acta Ecologica Sinica}, 2020.
\end{enumerate}

\section*{人工智能使用声明}

本文未使用任何人工智能工具。

\end{document}
